\scanpage{52}
\mspage{1}

I) Caractérisation de $W_\infty$ \hfill 19.11.90

Je voudrais caractériser $W_\infty$ en tant que partie de $\mathrm{Fl}(\mathrm{Cat})$, par ses propriétés comme localiseur, et sans faire appel à la connaissance de notions \unsure{utilisant|utilisées} \unsure{déjà|déf.} $W_\infty$, comme le proposait \unsure{Artin|etc.}. Je voudrais tout au moins des propriétés sur un localiseur $W$, familières pour $W_\infty$, qui \inserted{soient si \unsure{fortes} qu'elles} impliquent $W\subset W_\infty$. Bien sûr, dans les démonstrations je me permet d'utiliser tout ce que je sais sur $W_\infty$.

Bien sûr, il faudra au moins une forme de saturation, peut-être très faible, je présume que \unsure{suivant} suffit (non, j'ai dû \inserted{\unsure{le renforcer en (1)}})

% slanted
\margin{$W(1)$ Axiome de saturation \unsure{Ens}}

\textcircled{1} a) Les \deleted{\unsure{isomorphismes}}\inserted{identités} sont dans $W$ \deleted{\unsure{et de}} b) si $f,g$ composables et deux parmi $f,g,gf$ sont dans $W$, la troisième aussi, c) si $fg\in W$, $gf\in W$, alors $f,g\in W$.

% slanted
\margin{\uline{NB}}

Une \deleted{\unsure{axiome}} propriété particulièrement cruciale, qui déjà distingue $W_\infty$ de l'ens des homotopismes (notion trop stricte à beaucoup d'égards, est l'axiome que j'ai appelé L4 bis :

% slanted
\margin{$W(2)$ Axiome de localisation}

\textcircled{2} Soit $f:X\unsure{\text{---}|\longrightarrow}Y$ une flèche dans $\mathrm{Cat}_{/S}$, telle que $\forall s\in\mathrm{Ob}\,S$ $f_s:X_{/s}\longrightarrow Y_{/s}$ soit dans $W$ (on dira que $f$ est \deleted{\unsure{loc}} dans $W$ ``localement sur $S$''. Alors $f\in W$.

\deleted{\uline{Proposition}} \deleted{\unsure{Démon}} Le cas $Y=S$, cela donne

\uline{Corollaire} : Si $f:X\longrightarrow Y$ est $W$-asphérique i.e. $y\backslash X\longrightarrow y\backslash Y$ \deleted{\unsure{dans}} dans $W$ pour tout $y\in\mathrm{Ob}\,Y$, alors $f\in W$.

% slanted
\margin{On peut considérer \deleted{\unsure{cette}} \unsure{propriété} du cor. comme une version affaiblie $W(2\,\mathrm{bis})$ de $W(2)$ ci-dessus. Dans la \unsure{prop. 1} on se \unsure{borne} à $W(2\,\mathrm{bis})$, \unsure{au lieu de} $W(2)$ \deleted{\unsure{il y a dans ce cas a) b) c) sont équivalentes}}}

\uline{Proposition 1} Si $W$ satisfait \textcircled{1} et \textcircled{2}, \deleted{\unsure{deux}} \deleted{\unsure{assertions}} les propriétés suivantes \deleted{\unsure{sont équivalentes}} sont équivalentes :

a) Pour tout $X$ dans $\mathrm{Cat}$, $X\times\Delta^1\mathrel{\mathop{\longrightarrow}\limits^{\mathrm{pr}_1^X}}X$ est dans $W$ (i.e. $\Delta^1\longrightarrow e$ ``universellement dans $W$'').

b) Si $X$ dans $\mathrm{Cat}$ a un objet final, alors $X$ est $W$-asph. i.e. $X\longrightarrow e$ est dans $W$.

c) Tout homotopisme est dans $W$. \deleted{Alors $b\Rightarrow a,\ c\Rightarrow a$, et si les \unsure{flèches} sont \unsure{prises} \unsure{ensemble}, $a\Rightarrow b$, $a\Rightarrow c$, \unsure{b}.}

\uline{Dém} b) $\Rightarrow$ a), car pour voir que $\mathrm{pr}_1^X\in W$, en \unsure{appliquant} le corollaire ci-dessus, il faut prouver que $\forall x\in X$

\[
(X\times\Delta^1)_{/x}\longrightarrow X_{/x}\tag*{(\unsure{*|x})}
\]

est dans $W$. Or
\[
(X\times\Delta^1)_{/x}\simeq X_{/x}\times\Delta^1,
\]

\diagram{\begin{tikzcd}[ampersand replacement=\&,column sep=small,row sep=small]
\Delta^1 \arrow[d,dash] \& \deleted{\unsure{X}}X\times\Delta^1 \arrow[l] \arrow[d] \& X_{/x}\times\Delta^1 \arrow[l] \\
e \& X \arrow[l] \& X_{/x} \arrow[l,dash]
\end{tikzcd}}

et les deux catégories $X_{/x}$, $X_{/x}\times\Delta^1$ sont $W$-asph. car elles ont un objt final, donc on est \unsure{dans} $W$ par saturation \textcircled{1} : \deleted{car $X\times\Delta^1\longrightarrow X$ est \unsure{un homotopisme}.}

% slanted
\margin{on utilise \textcircled{1} (en sa forme faible, sans la condition de saturation c)}

% slanted
\margin{NB Seul $b\Rightarrow a$ utilise $W(2\,\mathrm{bis})$, les deux implications $a\Rightarrow c\Rightarrow b$ sont valables sous la seule hyp\textsuperscript{se} \textcircled{1} ; \deleted{$\unsure{c}\Rightarrow\unsure{a}$} $a\Rightarrow c$.}

Je rappelle que
\[
f:X\longrightarrow Y
\]

% archive sheet number: 52
